\documentclass[11pt]{article}
\usepackage[margin=1in]{geometry}
\usepackage{hyperref}
\usepackage{enumitem}
\usepackage{booktabs}
\usepackage{xcolor}
\title{ProbOS --- Month 3 Plan}
\author{Nisong Monyimba}
\date{\today}

\begin{document}
\maketitle

\section*{Status}
Per \texttt{docs/monthly\_plans/overall/main.tex}, Month 3 was always
DIRECTIONAL rather than pre-specified in detail -- Month 1's
retrospective explicitly flagged over-specifying future months as a
mistake, corrected by deferring detailed week-by-week planning until
the month is actually about to begin. This document intentionally
states THEMES, not a committed day-by-day schedule; the latter will
be written (following the same plan-then-retrospective process used
for every week so far) once Month 3's actual work begins.

\section*{Candidate themes, in rough priority order}

\subsection*{GPU kernel path}
Referenced as a genuine open item in
\texttt{docs/sbir\_phase1\_onepager.md}'s SBIR Ask (``GPU kernel, 10x
speedup over CPU, building on the existing C++/OpenMP path''). The
Month 1 Week 4 C++/OpenMP kernel achieved a measured 7x speedup over
the Python engine; a GPU path would need its own honest profiling
discipline (the same ``measure, don't assume'' standard that found
DP's real 1040x/45x wins alongside one genuine non-win in Month 1
Week 4) before any speedup number is claimed.

\subsection*{Extend the REST API's model registry}
Also an open SBIR ask: \texttt{python/server/main.py} currently only
registers \texttt{BatteryModel2Cell}. Extending it to the three Week 4
cross-discipline models (option pricer, ED queue, clinical trial) is
low-risk, since the pattern is already proven for one model and the
models themselves are unchanged.

\subsection*{PDSL v0.2 -- control flow}
Month 1 Week 4's PDSL compiler v0.1 supports only flat drift/diffusion
declarations, no loops or conditionals. Genuinely useful, but not
urgent -- no current model has needed it yet.

\subsection*{Apply the particle filter to real (non-synthetic) data}
The most consequential open SBIR ask (``apply the particle filter to
real sensor data from a pilot partner''), and the one most dependent
on external factors (an actual pilot partnership) rather than purely
technical work. Cannot be scheduled with certainty; flagged here so it
is not forgotten.

\subsection*{Continued quality-standard maintenance}
Per the mid-Month-2 hardening pass: \texttt{pre\_week\_audit.sh} and
\texttt{docs/standards/quality\_standards.md} are living documents.
Any new gap discovered (declared-but-unenforced checks, stale
documentation, vestigial code patterns) should be closed in the same
session it is found, not deferred -- the standing practice established
this month, to be continued.

\section*{Process}
As with every prior month, Month 3's actual week-by-week plan will be
written and compiled to PDF (\texttt{docs/monthly\_plans/month3/weekN/main.tex})
before that week's work begins, with a retrospective appended once it
completes. This document will be revisited and made more concrete once
Month 3 actually starts.

\end{document}
